\documentclass{thecours}
\begin{document}
    \bigpart{Définition et notations}
    \definition{equation}
    \medskip\par\noindent{\bfseries Exemple :}
    L'égalité suivante est une équation~:$$13+?=28$$
    Le nombre 15 est une solution de l'équation.%
    \medskip\par\noindent{\bfseries Exemple :}
    On préfère noter la variable avec une lettre (souvent $x$):$$13+x=28$$
    Cette équation possède une seule solution.
    \definition{resoudre_une_equation}
    \medskip\par\noindent{\bfseries Exemple :}
    L'équation suivante a 2 solutions (1 et 3)~:$$x^2=4x-3$$
    L'équation suivante n'a aucune solution~:$$x+1=x-1$$
    Tous les nombres sont solutions de l'équation suivante~:$$2(x+4)-3x=8-x$$
    \newpage\regle{resolution_d_une_equation}
    \methode{resoudre_une_equation}
    \bigpart{Équations du type $ax+b=cx+d$}
        \subpart{Équation du type $ax=d$}
        \subpart{Équation du type $ax+b=d$}
        \subpart{Équation du type $ax+b=cx+d$}
        Avant de résoudre, on peut~:
        développer, multiplier par le PPCM des dénominateurs dans le cas de fractions,
        réduire chaque membre
    \bigpart{Équations produit nul}
        \subpart{Produit nul}    
            On se pose la question suivante : comment puis-je choisir deux nombres $a$ et
            $b$ pour avoir $a\times b=0$~?
            Si $a$ et $b$ sont tous les deux différents de zéro, c'est impossible, il faut
            donc avoir $a=0$ ou $b=0$.
            \proposition{produit_nul}
            {\bfseries Exemple :}\par\noindent
            Résolution de l'équation $(x-3)(x+7)=0$.
        \subpart{Résolution de l'équation $x^2=a$}
            \proposition{equation_carre_constant}
            {\bfseries Exemple :}\par\noindent
            Résolution de l'équation $x^2=36$.
\end{document}